\documentclass[11pt]{article}

\usepackage[margin=1in]{geometry}
\usepackage[T1]{fontenc}
% Libertine text/math when available.
\IfFileExists{libertine.sty}{%
  \usepackage{libertine}%
  \IfFileExists{newtxmath.sty}{\usepackage[libertine]{newtxmath}}{}%
}{}
% Light monospace for inline code.
\IfFileExists{inconsolata.sty}{\usepackage[scaled=0.95]{inconsolata}}{%
  \IfFileExists{zi4.sty}{\usepackage[varqu,scaled=0.95]{zi4}}{}}
\usepackage{booktabs}
\usepackage{graphicx}
\usepackage{amsmath}
\usepackage{array}
\usepackage[table]{xcolor}
\usepackage[hidelinks]{hyperref}
\usepackage{xurl}
\usepackage{caption}

\captionsetup{font=small,labelfont=bf,skip=4pt}
\definecolor{hdrgray}{gray}{0.92}
\definecolor{keyrow}{RGB}{226,239,246}
\newcommand{\code}[1]{\texttt{#1}}
\newcommand{\best}[1]{\textbf{#1}}
\newcommand{\hdr}{\rowcolor{hdrgray}}
\newcommand{\key}{\rowcolor{keyrow}}
\setlength{\emergencystretch}{3em}

\newenvironment{ctab}{\par\medskip\centering\small
  \setlength{\tabcolsep}{5pt}\renewcommand{\arraystretch}{1.12}}{\par\medskip}

% Full-width abstract aligned with the body margins.
\renewenvironment{abstract}{%
  \par\medskip\noindent{\bfseries\abstractname.}\enspace\ignorespaces
}{\par\medskip}

\title{\textbf{The \code{get\_subtree} Tail on MongoDB Community Edition:\\
A Measured Optimization Study for the PageIndex Retrieval Workload}}
\author{Junyao Dong}
\date{}

\begin{document}
\maketitle

\begin{abstract}
A companion benchmark first reported MongoDB's \code{get\_subtree} P95 on a
ten-million-node tree at 2.5\,s, an order of magnitude above PostgreSQL, DuckDB,
and SQLite, and framed it as MongoDB's one weak spot on the PageIndex retrieval
workload. This report set out to turn that finding into an optimization plan
and then, because the benchmark harness was already in place, to \emph{measure}
the plan rather than argue it from documentation alone. Two results reshaped the
conclusion. First, the 2.5\,s figure does not reproduce: re-measured in
isolation and under concurrent load, with the collection resident in cache,
\code{get\_subtree} holds a stable P95 of roughly $0.3$--$0.4$\,s---about
$3\times$ PostgreSQL, not the $20\times$ the first run implied. The original tail
was an artifact of whole-system contention during the multi-engine run. Second,
the intuitive fix does not behave the way documented mechanism predicts.
Splitting the large \code{text} field into its own collection (the Subset
Pattern), the move we expected to matter most, leaves the P95 unchanged, because
on a host whose cache holds the whole collection the scan is gated by the
\emph{number} of per-document fetches, not their size. The one lever that helps
is the covering index we had initially demoted: served index-only it removes the
fetch and buys ${\sim}1.4\times$, though extending it to cover the multi-kilobyte
\code{summary} key inflates it to the size of the collection's own data. The
corrected picture is that \code{get\_subtree} is a sub-second,
${\sim}3\times$-PostgreSQL operation with a narrow optimization space; a larger
win needs structural denormalization (Nested Sets, a tree-order \code{\_id}, or a
pre-computed subtree document), which the write-once shape of the tree makes
viable but which we did not measure. Everything here stays on Community Edition,
no engine change.
\end{abstract}

\section{The problem, as first reported}

The PageIndex retriever reaches storage through four calls: a point lookup, an
expand-children, a content fetch, and \code{get\_subtree}, which pulls a node's
descendants to render the tree view and hand a block to the language model. On
the ten-million-node dataset the first three are sub-millisecond on MongoDB and
need no work. \code{get\_subtree} was the exception. A large subtree matches
about 36{,}000 nodes, and there the initial multi-engine run posted a MongoDB P50
of 30.9\,ms and a P95 of 2{,}478\,ms, against PostgreSQL at 10.7/123\,ms, SQLite
at 11.8/135\,ms, and DuckDB at 9.9/47\,ms (Table~\ref{tab:gap}). The median was
within a small factor of the others; the tail read more than an order of
magnitude worse, and that tail became the optimization target.

\begin{table}[h]
\caption{\code{get\_subtree} latency as first recorded, large dataset
(${\sim}$36k-node subtree, ms). The P95 shown here is the figure this report
re-examines.}
\label{tab:gap}
\begin{ctab}
\begin{tabular}{@{}lrrrr@{}}
\toprule
\hdr & MongoDB & PostgreSQL & DuckDB & SQLite \\
\midrule
P50 & 30.9  & 10.7 & \best{9.9} & 11.8 \\
\key P95 & 2{,}478 & 123  & \best{47} & 135 \\
\bottomrule
\end{tabular}
\end{ctab}
\end{table}

Two things then changed the story. The 2{,}478\,ms did not survive a controlled
re-measurement (Section~\ref{sec:remeasure}), and the fixes, once run rather than
reasoned about, did not rank the way the documentation suggested
(Section~\ref{sec:plan}). This report is the corrected account.

\section{Re-measuring the tail}\label{sec:remeasure}

The 2{,}478\,ms was recorded while the other three engines, the write benchmark,
and the concurrency sweep ran on the same host. To separate the query's own cost
from that contention we re-ran \code{get\_subtree} on the ten-million-node tree
under controlled conditions: one measuring client, the collection freshly loaded
and fully resident, the same 200 seeded subtree paths, optionally with $N$
background processes hammering point lookups at the same \code{mongod}.

In isolation the P95 lands at ${\sim}0.30$--$0.42$\,s (P50 ${\sim}30$--$40$\,ms)
and is stable across repeated runs---an order of magnitude below the first
figure. Driving the query while 16, 32, then 64 processes contend for the same
server barely moves it (Table~\ref{tab:remeasure}): P95 rises from 331 to
376\,ms, $1.1\times$, never approaching seconds. Throughout, the WiredTiger cache
holds the whole collection (${\sim}17$\,GB resident) and reports zero
application-thread evictions, so the tail is neither eviction nor ordinary read
contention.

\begin{table}[h]
\caption{\code{get\_subtree} re-measured on the 10M tree (one measuring client,
id-only projection, ms) while $N$ background processes drive concurrent point
lookups.}
\label{tab:remeasure}
\begin{ctab}
\begin{tabular}{@{}rrrrr@{}}
\toprule
\hdr Concurrent load $N$ & P50 & P95 & P99 & WT evictions \\
\midrule
0  & 31.4 & 331 & 476 & 0 \\
16 & 39.4 & 357 & 555 & 0 \\
32 & 44.1 & 415 & 534 & 0 \\
64 & 39.0 & 376 & 627 & 0 \\
\bottomrule
\end{tabular}
\end{ctab}
\end{table}

The corrected magnitude is milder than the first run implied: steady-state,
\code{get\_subtree} is ${\sim}0.4$\,s at P95 and ${\sim}30$--$40$\,ms at the
median, roughly $3\times$ PostgreSQL at \emph{both} (its 123\,ms tail and 11\,ms
median), not the ${\sim}20\times$ the original P95 alone suggested. The 2{,}478\,ms
was a property of the loaded machine---most likely whole-system resource
contention, or a cold-cache first touch before the collection was warm---not of
the query in steady state. The mechanism below still holds; its magnitude is
sub-second.

\section{Root cause: the count of fetches, not their size}

A secondary index in MongoDB (here on \code{path}) is a separate B-tree holding
only the indexed value and a record id. To return any other field, WiredTiger
must \emph{fetch} the full BSON document and apply the projection to it. A point
lookup is one such fetch and is cheap. \code{get\_subtree} is tens of thousands
of them. The data is cache-resident---and pages in the WiredTiger cache are held
\emph{uncompressed}---so the cost is neither disk I/O nor decompression. It is the
per-document work of locating each record, materializing the document, applying
the projection, and streaming the result through the cursor, multiplied by
subtree size.

The initial hypothesis pinned that per-document cost on the wide leaf
\code{text}: each fetch materializes the whole document, \code{text} payload and
all, only to project it away, so removing \code{text} should make each fetch
cheap. The measurement in the next section refutes that reading. Shrinking each
document by ${\sim}75\%$ leaves \code{get\_subtree} unchanged, because on this host
document \emph{size} never gates the scan---document \emph{count} does. The cost
is the ${\sim}36$k separate \code{FETCH} operations, and each one is dominated by
locating the record and threading it through the cursor, not by the number of
bytes in it once it is already resident and uncompressed in cache. The
consequence for optimization is sharp: only \emph{removing} fetches helps.
Shrinking them does not, and neither does changing the query operator.

This is also why the relational engines win the operation specifically: a row
store visits each matching row but materializes only the requested
column---PostgreSQL does not even decompress the JSONB datum, since
\code{node\_id} is a separate column---and a column store reads only the requested
column in vectorized batches. They do less work per node, and there are tens of
thousands of nodes. The same document model that costs here is what makes MongoDB
fast everywhere else: storing a node's fields together makes ``fetch one whole
entity'' cheap (point and children reads), makes in-place updates cheap (no vacuum
debt), and keeps the on-disk form compact. The workload plays to that model on
every operation but this one.

\section{The optimization plan, measured}\label{sec:plan}

Because the harness was already in place, we did not stop at recommending fixes;
we ran them on the ten-million-node tree and measured (Table~\ref{tab:opt}). Each
stays on Community Edition with no engine change. The results invert the intuitive
ranking: the Subset Pattern, the move we expected to lead, does nothing here, and
the covering index we had demoted is the only lever that moves the number.

\begin{table}[h]
\caption{\code{get\_subtree} fixes measured on the 10M tree (one client). ``view''
returns \code{node\_id}+\code{title}+\code{summary} (what the tree view renders);
``id'' returns \code{node\_id} only. P50/P95 in ms.}
\label{tab:opt}
\begin{ctab}
\begin{tabular}{@{}p{6.0cm}rrl@{}}
\toprule
\hdr Variant & P50 & P95 & Result \\
\midrule
\key Baseline, view (\code{text} inline) & 39.6 & 419 & reference \\
Subset Pattern (split \code{text}), view & 39.2 & 419 & \best{no change} \\
Lean covering \code{\{path,node\_id\}}, id & 19.3 & 191 & ${\sim}1.4\times$, covered \\
Fat covering \code{+title,summary}, view & 33.5 & 350 & ${\sim}1.2\times$; 4.66\,GB index \\
Lean covering ${+}$ batched \code{\$in}, view & 137 & 1314 & ${\sim}3\times$ slower \\
\bottomrule
\end{tabular}
\end{ctab}
\end{table}

\subsection{Splitting the \code{text} field does nothing (here)}

The single most intuitive move---the Subset Pattern, splitting \code{text} into
its own collection keyed by node id so the scanned documents are small---was the
leading hypothesis, and on this host it is wrong. It shrinks the collection by
${\sim}75\%$ ($4.97\to1.25$\,GB of data) yet leaves \code{get\_subtree} P95
exactly where it was ($419\to419$\,ms for the view, $301\to284$\,ms for id-only,
within run-to-run noise). The reason is the one from the previous section: the
cache holds the whole collection either way, so document size never gates the
scan and the ${\sim}36$k fetches remain ${\sim}36$k fetches.

The design still costs the hot path nothing---\code{get\_subtree} never needs
\code{text}, and the body is read later and separately by \code{get\_entity}, on
the few nodes the model selects---so it remains a reasonable schema choice on its
own merits (a smaller working set, a lighter write path). What it is \emph{not},
on this workload and this host, is a \code{get\_subtree} fix. \emph{Caveat, and it
is the important one:} the Subset Pattern's entire premise is a working set too
large for RAM, and this experiment ran with a WiredTiger cache many times the
collection size, so it could not have exercised that premise. On a
cache-constrained production host, where the collection with \code{text} spills
RAM and the version without it fits, the split could matter a great deal. Our
result says only that document \emph{size} is not the lever \emph{when the data is
resident}; it does not say the split is worthless in general. This is the one open
retest worth running (Section~\ref{sec:open}).

\subsection{A covering index is the only real lever, and a modest one}

Removing the \code{FETCH} is the one thing that helps. A lean
\code{\{path,node\_id\}} index answers the id-only subtree query index-only---
\code{explain} confirms \code{PROJECTION\_COVERED}, \code{totalDocsExamined:\,0},
\code{keysExamined:\,12838}---for ${\sim}1.4\times$ (P95 $301\to191$\,ms). This is
the fix the earlier hypothesis demoted, on the theory that the multi-kilobyte
\code{summary} key would defeat a covering index; the id-only index carries no
\code{summary} and is not defeated.

The complication is that the tree view needs \code{title} and \code{summary}, not
just \code{node\_id}. Extending the index to cover them
(\code{\{path,node\_id,title,summary\}}) does work---it is still a fully covered
scan, \code{totalDocsExamined:\,0}, and it beats the baseline view
($419\to350$\,ms, ${\sim}1.2\times$)---but it inflates the index to 4.66\,GB, as
large as the collection's own data, because \code{summary} really is a
multi-kilobyte key and WiredTiger's index prefix compression deduplicates the
shared \code{path} prefixes but does nothing for the fat \code{summary} values. So
the fat-key concern was real: the covering index helps, but the cheap version
($1.4\times$) covers only the id projection, and covering the actual view
projection costs an index the size of the data for a $1.2\times$ gain.

Resolving the view a third way---a covered id scan followed by a batched
\code{\$in} to fetch \code{title} and \code{summary}---is worse, not better
(${\sim}3\times$ slower, P95 1{,}314\,ms). A ${\sim}36$k-id \code{\$in} overflows
the 16\,MB command limit, forcing ${\sim}36$ chunked round trips, and each chunk
re-fetches the same documents the plain scan would have, with cursor overhead on
top. Keep the anchored materialized-path range as the predicate---an unanchored
regex cannot form a range and must scan the whole index---and treat cache sizing
and \code{zstd} as irrelevant here: the data is already resident and held
uncompressed, so neither touches the per-fetch cost.

\subsection{A larger win needs structural denormalization (untested)}

The covering index shaves fetches at the margin; the structural fixes remove them
wholesale, and are the only route to closing the remaining ${\sim}3\times$ to
PostgreSQL. All three below are viable because the tree is largely write-once
after ingest, so their cost is paid once, at load, not on every read.

\paragraph{Nested Sets.} Assign each node a \code{left}/\code{right} number in a
depth-first traversal at ingest; a subtree is then a single range query,
\code{find(\{left:\{\$gt:p.left\}, right:\{\$lt:p.right\}\})}, with no recursion
and no per-node fetch beyond the range. MongoDB recommends this for static trees
and warns it is ``inefficient for modifying the tree structure'' because inserts
and deletes require $O(n)$ renumbering. The workload's incremental-insert path
(reindexing adds nodes after ingest) means the tree is not \emph{strictly}
write-once, so this fits a load-then-serve deployment or a batched re-number at
reindex time, not a steadily growing tree renumbered per insert.

\paragraph{Tree-order \code{\_id} on a clustered collection.} A clustered
collection stores documents physically ordered by \code{\_id}, in the same file as
the index, so a subtree range scan becomes a sequential read rather than a series
of random fetches---provided tree order (a DFS or nested-set number) is encoded
into \code{\_id}. The optimizer will not auto-select the clustered index when a
usable secondary index exists, so the query needs a hint.

\paragraph{Pre-computed subtree document.} Because the tree is write-once in the
common case, the structure-and-summary view of each bounded subtree root can be
materialized once at ingest and stored as a single document, turning
\code{get\_subtree} into a point read. It must exclude \code{text} to fit the
16\,MiB document limit ($1.5$\,kB\,${\times}$\,36k\,${\approx}55$\,MB with
\code{text}; structure-and-summary buckets of a few levels stay well under it).
The open cost is storage amplification from overlapping subtrees.

We did not measure any of these. With the baseline already at ${\sim}0.4$\,s on an
infrequent operation, the honest open question is whether the schema complexity is
worth shaving a sub-second tail---which is a product decision, not a mechanism
one. (\code{\$graphLookup} is not a candidate: it re-traverses \code{parent\_id},
still materializes documents, and in 7.0 is capped at 100\,MB and ignores
\code{allowDiskUse}.)

\section{A newer MongoDB version does not close the gap}

Upgrading the engine does not fix this without a schema change. Against this
operation---a projection-heavy range scan over ordinary documents---the 8.0--8.3
line offers nothing:

\begin{itemize}\itemsep2pt
\item The 8.0 \textbf{Express execution path} is a fast path for point and
equality lookups on a single index (\code{EXPRESS\_IXSCAN},
\code{EXPRESS\_CLUSTERED\_IXSCAN}): \code{\_id}-style queries with no range, sort,
or skip. It does not apply to a materialized-path range scan.
\item 8.0 \textbf{block (vectorized) processing} is scoped to time-series
collections and does not extend to general collections through 8.3.
\item The \textbf{\code{\$graphLookup} disk-spill fix} (SERVER-23980, lifting the
100\,MB cap that ignores \code{allowDiskUse}) lands only in 8.2.0---reverted from
8.1---and even then only helps the recursive-traversal alternative, which loses to
the indexed range scan here anyway.
\end{itemize}

\noindent The document-at-a-time \code{FETCH} is unchanged across the line, so the
remedy is schema and index design, not a version bump.

\section{What Community Edition does not put on the table}

The feature that would solve this most directly is not available without Atlas or
Enterprise. There is no \textbf{column-store index} (Atlas-only)---the one feature
that would make a projection-heavy scan read only the projected column, the way
the relational engines do. But PageIndex retrieval never scans the whole corpus,
so a column store solves a problem this workload does not have; the schema fixes
above address the one it does. Self-managed \textbf{Search and Vector Search}
became runnable on Community in 8.2 (a separate \code{mongot} process, public
preview, replica set required), but they serve relevance retrieval, not a
range-projection scan, so they do not address \code{get\_subtree} either.
Auto-scaling, stream processing, tiered storage, and federated queries are
Atlas-only; the in-memory engine and auditing are Enterprise-only. \emph{No Atlas-
or Enterprise-only feature would solve \code{get\_subtree}}: the operation is bound
by the document-at-a-time read model, and the only levers are schema, index, and
storage configuration---all present on Community.

\section{What remains open}\label{sec:open}

Two things the measurements did not settle. First, the exact origin of the
original 2{,}478\,ms: point-lookup contention alone does not reproduce it
(Section~\ref{sec:remeasure}), so the cause was heavier than the concurrency the
workload itself generates---whole-system resource pressure during the multi-engine
run, or a cold-cache first touch. Second, and more useful to close: the Subset
Pattern under a \emph{cache-constrained} host. This study ran with the collection
fully resident, which by construction cannot show the split's benefit; a re-run
with \code{--wiredTigerCacheSizeGB} set below the working set (collection with
\code{text} spilling RAM, without it fitting) would test the premise the pattern
actually targets. And the measured payoff of the structural-denormalization
options remains unquantified. None of these changes the practical conclusion:
\code{get\_subtree} is a sub-second, ${\sim}3\times$-PostgreSQL operation whose one
cheap fix---a covering index---has limited upside.

\section{Conclusion}

MongoDB's one weakness on the PageIndex workload, the \code{get\_subtree} tail, is
real but smaller than first reported. A controlled re-measurement puts its
steady-state P95 at ${\sim}0.4$\,s, roughly $3\times$ PostgreSQL and ${\sim}3\times$
at the median too---not the multi-second tail the first run recorded, which traces
to contention rather than to the query. The cost is the \emph{count} of per-node
fetches, which is why the intuitive fix fails and the demoted one wins: splitting
\code{text} shrinks each document but not their number and leaves the P95
unchanged on a resident collection, while a covering index removes the fetch
outright and is the only cheap lever---worth just ${\sim}1.4\times$, and only
${\sim}1.2\times$ once it must carry the fat \code{summary} key. A larger win needs
structural denormalization the write-once tree makes viable but that we did not
measure. The write, concurrency, and footprint advantages that make MongoDB
attractive elsewhere are untouched by any of this. The remaining work is not more
analysis---it is the cache-constrained Subset retest and the one structural
before/after that would price the larger win.

\section*{References}
\small
\begin{enumerate}\itemsep2pt
\item Subset Pattern. \url{https://www.mongodb.com/docs/manual/data-modeling/design-patterns/group-data/subset-pattern/}
\item Model tree structures with Nested Sets. \url{https://www.mongodb.com/docs/v7.0/tutorial/model-tree-structures-with-nested-sets/}
\item Model tree structures with Materialized Paths. \url{https://www.mongodb.com/docs/manual/tutorial/model-tree-structures-with-materialized-paths/}
\item Clustered collections. \url{https://www.mongodb.com/docs/v7.0/core/clustered-collections/}
\item WiredTiger storage engine (cache holds data uncompressed; cache sizing). \url{https://www.mongodb.com/docs/manual/core/wiredtiger/}
\item Query optimization / covered queries; explain results (covered = IXSCAN with no FETCH). \url{https://www.mongodb.com/docs/manual/reference/explain-results/}
\item \code{\$graphLookup} memory limit. \url{https://www.mongodb.com/docs/v7.0/reference/operator/aggregation/graphlookup/}
\item MongoDB 8.0 release notes (Express path; block processing). \url{https://www.mongodb.com/docs/manual/release-notes/8.0/}
\item SERVER-23980 (\code{\$graphLookup} disk spill, fixed 8.2.0); SERVER-102453 (reverted from 8.1). \url{https://jira.mongodb.org/browse/SERVER-23980}
\item MongoDB Community Edition feature scope (column-store and analytics are Atlas/Enterprise). \url{https://www.mongodb.com/products/self-managed/community-edition}
\item Independent PostgreSQL-vs-MongoDB benchmark, off-workload (direction only). \url{https://www.mdpi.com/2504-2289/10/2/66}
\end{enumerate}

\end{document}
